% page 219 du fac-simile (image p-218 du scan)
% bloc 167.6 x 233.3 mm ; marge gauche 34.4 mm
\pgc{12.700mm}{0.000mm}{
\l{\cel{54}211}
\l{}
\l{\sou{helixo}. (\sou{anat.}) La rebordo di la orel-funelo. - EFIS.}
\l{}
\l{\sou{\textquotedbl{}heller\textquotedbl{}}. (\sou{ante la milito universala di 1914-18}). Moneto-peco}
\l{\cel{3}ek bronzo e nikelo, adoptita da Austria-Hungaria, centima}
\l{\cel{3}parto di la \textquotedbl{}krono\textquotedbl{}.}
\l{}
\l{helmo. I. (\sou{olim}). Kasko pintoza qua kovras la kapo e la vizajo,}
\l{\cel{3}kun aperturo por la okuli. -II. (\sou{blazon-arto}). Peco qua in-}
\l{\cel{3}dikas, per sua situeso, la gradi diversa di nobeleso,\cel{2}DEFIS.}
\l{}
\l{helpar. (\sou{trans.}) Partoprenar lo facata od agata da ulu qua ne}
\l{\cel{3}povas facar od agar lo sole. - DE.}
\l{}
\l{\sou{hem}\cel{1}!\cel{2}Interjeciono qua expresas dubito, hezito.}
\l{}
\l{\sou{hemo}. La loko intima ube onu vivas (kontraste kun loko ube onu}
\l{\cel{3}vivas, en hotelo o lojeyo). - DEF.}
\l{}
\l{\sou{hematito}. (\sou{mineral.}) Sanguino, erco reda brunatre.-\cel{2}DEFIRS.}
\l{}
\l{\sou{hemerokalo.}\cel{2}(\sou{bot.}) Planto liliacea, di qua la speco precipua}
\l{\cel{3}esas la konvolvulo. - L.\cel{1}\sou{Hemerocallis}. - EFIS.}
\l{}
\l{\sou{hemiopa}. (\sou{patol.)}\cel{2}Malade de ta perturbeso vidala, pro qua onu di-}
\l{\cel{3}cernas nur la duimo sinistra o dextra di la objekto. - EF.}
\l{}
\l{hemiplegio. (\sou{patol.}) Paralizeso di duimo laterala di la homo-}
\l{\cel{3}korpo. - DEF.}
\l{}
\l{hemiptero. (\sou{zool.}) Insekto karakterizata da mi-elitri qui kovras}
\l{\cel{3}la parto avana di la ali. - DEFIS.}
\l{}
\l{hemistiko. (\sou{prosodio}). Duimo de verso, separita de la altra duimo}
\l{\cel{3}per silenceto (cezuro). - DEFIS.}
\l{}
\l{hemoglobino. (kemio biol.). Pigmento nitroza e feroza, reda,quan}
\l{\cel{3}onu renkontras en la hematii di la vertebrozi ed en la hemo-}
\l{\cel{3}limfo di ula senvertebri. - DEF.}
\l{}
\l{hemolizo. (\sou{biol.}) Destrukto di la sango-globuli reda.}
\l{}
\l{hemoptizio. (\sou{patol}.) Sputado di sango. - DEF.}
\l{}
\l{hemoragio. (\sou{patol.}) Fluo di sango ek la vaskuli, kun o sen rupt-}
\l{\cel{3}eso di la parieti. - DEF.}
\l{}
\l{hemoroido. (\sou{patol.)}\cel{2}Tumoro di la veini di la anuso, de qua sango}
\l{\cel{3}eskapas intervalope. - DEFIRS.}
\l{}
\l{\textquotedbl{}\sou{henry\textquotedbl{}.}\cel{2}(\sou{elektrotekniko).}\cel{1}Unajo di indukto reciproka o di auto-}
\l{\cel{3}indukto, di qua la valoro egalesas 109 g unaji elektromagnet-}
\l{\cel{3}ala C.G.S.}
}
